% Gemini theme
% https://github.com/anishathalye/gemini

\documentclass[final]{beamer}

% ====================
% Packages
% ====================

\usepackage[T1]{fontenc}
\usepackage{lmodern}
\usepackage[size=custom,width=75,height=100,scale=1.0]{beamerposter}
\usetheme{gemini}
\usecolortheme{cse19}
\usepackage{graphicx}
\usepackage{booktabs}
\usepackage{tikz}
\usepackage{pgfplots}
\pgfplotsset{compat=1.15}

% ====================
% Lengths
% ====================

% If you have N columns, choose \sepwidth and \colwidth such that
% (N+1)*\sepwidth + N*\colwidth = \paperwidth
\newlength{\sepwidth}
\newlength{\colwidth}
\setlength{\sepwidth}{0.0333\paperwidth}
\setlength{\colwidth}{0.45\paperwidth}

% Spacing / Formatting
\newcommand{\separatorcolumn}{\begin{column}{\sepwidth}\end{column}}
\newcommand{\itembox}[1]{\item {\makebox[7cm]{#1 \hfill}}}

% Standard Probability
\DeclareMathOperator{\E}{\mathbb{E}}
\newcommand{\PP}{\mathbb{P}}
\newcommand{\RR}{\mathbb{R}}

% Data-Consistent Inversion Framework

\DeclareMathOperator*{\param}{\lambda}
\DeclareMathOperator*{\pspace}{\Lambda}
\newcommand{\pmeas}{\mu_{\pspace}}
\newcommand{\pborel}{\mathcal{B}_{\pspace}}

\DeclareMathOperator*{\initialP}{\PP^{in}}
\DeclareMathOperator*{\initial}{\pi^{in}}
\DeclareMathOperator*{\updatedP}{\PP^{up}}
\DeclareMathOperator*{\updated}{\pi^{up}}

\DeclareMathOperator*{\noise}{\xi}
\DeclareMathOperator*{\nspace}{\Xi}

\newcommand{\obs}{\boldsymbol{o}}
\newcommand{\data}{\boldsymbol{d}}
\DeclareMathOperator*{\dspace}{\mathcal{D}}
\newcommand{\dmeas}{\mu_{\dspace}}
\newcommand{\dborel}{\mathcal{B}_{\dspace}}

\DeclareMathOperator*{\lam}{\left ( \param \right ) }
\DeclareMathOperator*{\qoi}{Q \lam }
\DeclareMathOperator*{\qlam}{\left ( \qoi \right ) }
\DeclareMathOperator*{\M}{\mathcal{M}}

\DeclareMathOperator*{\observedP}{\PP^{ob}}
\DeclareMathOperator*{\observed}{\pi^{ob}}
\DeclareMathOperator*{\predictedP}{\PP^{pr}}
\DeclareMathOperator*{\predicted}{\pi^{pr}}

\DeclareMathOperator*{\dciP}{\updatedP = \initialP \frac{\observedP}{\predictedP}}
\DeclareMathOperator*{\dciD}{\updated = \initial \frac{\observed}{\predicted}}
\DeclareMathOperator*{\dci}{\updated\lam = \initial\lam \frac{\observed\qlam}{\predicted\qlam}}

% ====================
% Title
% ====================

\title{Push-forward Measures for Parameter Identification under Uncertainty}

\author{Michael Pilosov \inst{1}}

\institute[shortinst]{\inst{1} University of Colorado: Denver}

% ====================
% Body
% ====================

\begin{document}

\begin{frame}[t]
\begin{columns}[t]
\separatorcolumn

\begin{column}{\colwidth}

  \begin{block}{Notation}
\large
    \begin{itemize}
        \itembox{$\pspace \subset \RR^P$} Parameter Space
        \itembox{$\obs: \pspace \to \mathcal{O} \subset \RR^D$} Observables
        \itembox{$\nspace \subset \RR^D$} Noise Space
        \itembox{$\param^\dagger\in\pspace$} True (Reference) Parameter
        \itembox{$\data(\noise) \subset \RR^D$} Possible Data, where $d_i(\noise) = \obs_i(\param^\dagger) + \xi_i$
        \itembox{$\noise^\dagger\in\nspace$} Noise that Polluted (Observed) Data
        \itembox{$\data^\dagger\in\RR^D$} Observed Data, where $\data^\dagger = \data(\noise^\dagger)$
    \end{itemize}

\Large
    \heading{Quantity of Interest Map}
    \emph{\Large Defines a (functional) relationship between predictions and data}

        $$Q \left (\param, \noise \right ) = F \left ( \obs(\param), \data(\noise) \right ) $$
        $$Q \left (\param \right ) = F \left ( \obs(\param), \data^\dagger \right )$$
        $$Q \left ( \pspace \right ) := \dspace \subset \RR$$
\normalsize
    \begin{figure}
        \includegraphics[width=20cm]{ref1}
    \end{figure}

\end{block}
\vspace{-1cm}


\end{column}

\separatorcolumn

\begin{column}{\colwidth}


  \begin{block}{Framework}
\large
    \begin{itemize}
        \itembox{$ \PP $} Probability Measure
        \itembox{$ \pi $} Density (Radon-Nikodym)
        \itembox{$ \initialP, \; \initial $} Initial
        \itembox{$ \observedP, \; \observed $} Observed
        \itembox{$ \predictedP, \; \predicted $} Predicted ({\bf Push-forward of Initial})
        \itembox{$ \updatedP, \; \updated $} Updated
        \itembox{$ \sigma $} Variance of Noise
    \end{itemize}

\Large
    \heading{Updating Initials with Observations and Predictions}
    \begin{equation*}
            \dciP 
    \end{equation*}
    \begin{equation*}
            \dci
    \end{equation*}

\end{block}
\vspace{-1cm}


  \begin{block}{Nullam vel erat at velit convallis laoreet}

    Class aptent taciti sociosqu ad litora torquent per conubia nostra, per
    inceptos himenaeos. Phasellus libero enim, gravida sed erat sit amet,
    scelerisque congue diam. Fusce dapibus dui ut augue pulvinar iaculis.

    \begin{table}
      \centering
      \begin{tabular}{l r r c}
        \toprule
        \textbf{First column} & \textbf{Second column} & \textbf{Third column} & \textbf{Fourth} \\
        \midrule
        Foo & 13.37 & 384,394 & $\alpha$ \\
        Bar & 2.17 & 1,392 & $\beta$ \\
        Baz & 3.14 & 83,742 & $\delta$ \\
        Qux & 7.59 & 974 & $\gamma$ \\
        \bottomrule
      \end{tabular}
      \caption{A table caption.}
    \end{table}

  \end{block}

  \begin{block}{References}
    \centering
    \begin{figure}
        \includegraphics{ref1}
    \end{figure}
   
    %\nocite{*}
    %\footnotesize{\bibliographystyle{plain}\bibliography{poster}}

  \end{block}

\end{column}

\separatorcolumn
\end{columns}
\end{frame}

\end{document}
